
%%%%%%%%%%%%%%%%%%%%%%%%%%%%%%%%%%%%%%%%%%%%%%%
%                                             %
% MOJE PAKIETY, USTAWIENIA ITD                %
%                                             %
%%%%%%%%%%%%%%%%%%%%%%%%%%%%%%%%%%%%%%%%%%%%%%%

% Tutaj proszę umieszczać swoje pakiety, makra, ustawienia itd.

% Pakiety dodatkowe (booktabs, graphicx, subcaption już są w settings.tex)
\usepackage{float}            % środowisko [H] - tabele/figury w miejscu
\usepackage{xcolor}           % kolory w listingach
\usepackage{tabularx}         % tabele z elastyczną szerokością kolumn
\usepackage{multirow}         % komórki łączone pionowo
\usepackage{makecell}         % wieloliniowe komórki tabel
\usepackage{microtype}        % lepszy kerning, hyphenation
% Pozwala używać środowiska "comment" do tymczasowego komentowania fragmentów
\usepackage{comment}

%%%%%%%%%%%%%%%%%%%%%%%%%%%%%%%%%%%%%%%%%%%%%%%%%%%%%%%%%%%%%%%%%%%%%
% listingi i fragmenty kodu źródłowego (Python - dla skryptów ML)
% % % % % % % % % % % % % % % % % % % % % % % % % % % % % % % % % % %

\usepackage{listings}

\definecolor{codegray}{rgb}{0.5,0.5,0.5}
\definecolor{codepurple}{rgb}{0.58,0,0.82}
\definecolor{codegreen}{rgb}{0,0.5,0}
\definecolor{codebg}{rgb}{0.97,0.97,0.97}

\lstdefinestyle{pythonstyle}{%
    language=Python,
    backgroundcolor=\color{codebg},
    commentstyle=\color{codegreen}\itshape,
    keywordstyle=\color{blue!70!black}\bfseries,
    numberstyle=\tiny\color{codegray},
    stringstyle=\color{codepurple},
    basicstyle=\ttfamily\footnotesize,
    breakatwhitespace=false,
    breaklines=true,
    captionpos=b,
    keepspaces=true,
    numbers=left,
    numbersep=8pt,
    showspaces=false,
    showstringspaces=false,
    showtabs=false,
    tabsize=4,
    frame=single,
    rulecolor=\color{black!20}
}

% Domyślny styl listingów - Python.
% (Szablon uczelni narzuca C++ jako domyślny - my pracujemy w Pythonie, nadpisujemy)
\lstset{style=pythonstyle}

%%%%%%%%%%%%%%%%%%%%%%%%%%%%%%%%%%%%%%%%%%%%%%%%%%%%%%%%%%%%%%%%%%%%%

% Własne komendy dla spójnego formatowania

% Skróty dla metryk - wszędzie w pracy formatowane jednakowo
\newcommand{\fone}{F\textsubscript{1}}
\newcommand{\precision}{\textit{precision}}
\newcommand{\recall}{\textit{recall}}
\newcommand{\accuracy}{\textit{accuracy}}

% Wyróżnienie nazw modeli i bibliotek
\newcommand{\xlmr}{XLM-RoBERTa}
\newcommand{\bert}{BERT}

%%%%%%%%%%%%%%%%%%%%%%%%%%%%%%%%%%%%%%%%%%%%%%%
%                                             %
% KONIEC MOICH USTAWIEŃ                       %
%                                             %
%%%%%%%%%%%%%%%%%%%%%%%%%%%%%%%%%%%%%%%%%%%%%%%
